% Generated by scripts/assemble.py; corrections belong in source/corrections.json.
\examhead{120}{南昌大学 2025 年数学分析试卷(回忆版)}{115}
% p116-b09; source PDF page 116
\par\Needspace{5\baselineskip}
\blockanchor{p116-b09}
\majorquestion{一}填空题. 每题 5 分, 共 30 分.\par
% p116-b10; source PDF page 116
\blockanchor{p116-b10}
\question{1}$\lim_{x\to 0}\frac{\int_0^x\sum_{n = 0}^\infty(-1)^n\frac{t^{2n + 1}}{2^n(2n + 1)!}\mathrm{d}t - \frac{x^2}{2}}{x(\sqrt[3]{1 + x} + e^x)}$\par
% p116-b11; source PDF page 116
\blockanchor{p116-b11}
\question{2}$\int \frac{\mathrm{d}x}{1 + \sqrt{1 - x^2}} =$\par
% p116-b12; source PDF page 116
\blockanchor{p116-b12}
\question{3}$\sum_{n=1}^{\infty} \ln \left( \cos \frac{1}{n^3} \right) x^n$ 收敛半径为\par
% p116-b13; source PDF page 116
\blockanchor{p116-b13}
\question{4}$x + y + z + 1 = 0$  到椭球  $18x^{2} + ()y^{2} + ()z^{2} = ()$  的距离 d = \_\_\_\_.\par
% p116-b14; source PDF page 116
\blockanchor{p116-b14}
\question{5}$u = x^{y} y^{z} z^{x}$  的全微分 du = \_\_\_\_.\par
% p116-b15; source PDF page 116
\blockanchor{p116-b15}
\question{6}$x^{2} + y^{2} = 6 - z, xz$ 平面, $yz$ 平面与 $y = 4z$, $x = 1$, $y = 2$ 围成体积 $V = \_$.\par
% p116-b16; source PDF page 116
\par\Needspace{5\baselineskip}
\blockanchor{p116-b16}
\majorquestion{二}(10 分)  $\lim_{n\to\infty}a_n=a,\quad|q|<1,$  用  $\varepsilon-N$  语言证明:  $\lim_{n\to\infty}(a_n+a_{n-1}q+\cdots+a_1q^{n-1})=\frac{a}{1-q}.$\par
% p116-b17; source PDF page 116
\par\Needspace{5\baselineskip}
\blockanchor{p116-b17}
\majorquestion{三}(10 分) $|f(x) - f(y)| \leq M|x - y|$ 为 Lipschitz 条件, $f(x)$ 满足 $L$ 条件, $a > 0$, 证明: $\frac{f(x)}{x}$ 在 $[a, +\infty)$ 上一致连续.\par
% p116-b18; source PDF page 116
\par\Needspace{5\baselineskip}
\blockanchor{p116-b18}
\majorquestion{四}(10 分) 设 $f(x)$ 二阶可导, 且 $|f(x)| \leq 1$, $|f''(x)| \leq 2$, 证明: $|f'(x)| \leq 3$. (区间可能是 [0, 2])\par
% p116-b19; source PDF page 116
\par\Needspace{5\baselineskip}
\blockanchor{p116-b19}
\majorquestion{五}(10 分) 设 $f(x)$ 在 $[a, b]$ 上无界, 证明: 对任意的 $\xi \in [a, b]$, 存在 $\delta > 0$, 使得 $f$ 在 $(\xi - \delta, \xi + \delta) \cap [a, b]$ 上无界.\par
% p116-b20; source PDF page 116
\par\Needspace{5\baselineskip}
\blockanchor{p116-b20}
\majorquestion{六}(10 分) 证明无穷积分的阿贝判别法, 即若无穷积分 $\int_{a}^{+\infty} f(x) \mathrm{d}x$ 收敛, $g(x)$ 在 $[a, +\infty)$ 上单调有界, 则 $\int_{a}^{+\infty} f(x) g(x) \mathrm{d}x$ 收敛.\par
% p116-b21; source PDF page 116
\par\Needspace{5\baselineskip}
\blockanchor{p116-b21}
\majorquestion{七}(10 分) 证明: $f(x) = \begin{cases} x^2, & [0,1] \text{ 所有有理数;} \\ 0, & [0,1] \text{ 所有无理数.} \end{cases}$ 不可积.\par
% p116-b22; source PDF page 116
\par\Needspace{5\baselineskip}
\blockanchor{p116-b22}
\majorquestion{八}(15 分) $f(x) = \begin{cases} \pi - x, & 0 < x \leq \pi; \\ 0, & x = 0; \\ -\pi - x, & -\pi \leq x < 0. \end{cases}$\par
% p117-b00; source PDF page 117
\blockanchor{p117-b00}
(1) 写出 $f(x)$ 的傅里叶级数.\par
% p117-b01; source PDF page 117
\blockanchor{p117-b01}
(2) 证明 $f(x)$ 在 $(- \pi, \pi)$ 上的一致收敛性(应该是傅里叶级数的一致收敛性).\par
% p117-b02; source PDF page 117
\par\Needspace{5\baselineskip}
\blockanchor{p117-b02}
\majorquestion{九}(15 分) 函数项级数 $\sum_{n=1}^{\infty} ne^{-nx}$.\par
% p117-b03; source PDF page 117
\blockanchor{p117-b03}
(1) 证明级数在 $(0, +\infty)$ 上收敛但不一致收敛.\par
% p117-b04; source PDF page 117
\blockanchor{p117-b04}
(2) 求级数的和函数.\par
% p117-b05; source PDF page 117
\par\Needspace{5\baselineskip}
\blockanchor{p117-b05}
\majorquestion{十}(15 分) 求 $I(\alpha) = \int_{0}^{+\infty} \frac{\arctan x}{x^{\alpha}(2 + x^{3})} \, \mathrm{d}x$ 的连续区间.\par
% p117-b06; source PDF page 117
\par\Needspace{5\baselineskip}
\par\Needspace{6\baselineskip}
\blockanchor{p117-b06}
\majorquestion{十一}(15 分) 计算\par
\nopagebreak[4]
% p117-b07; source PDF page 117
\blockanchor{p117-b07}
\displayformula{\iint_ {\Sigma} (x ^ {2} \cos (\boldsymbol {n}, x) + y ^ {2} \cos (\boldsymbol {n}, y) + z ^ {2} \cos (\boldsymbol {n}, z)) \mathrm{d} S.}
% p117-b08; source PDF page 117
\blockanchor{p117-b08}
其中 $S: x^{2} + y^{2} = z^{2} (0 \leq z \leq h)$, $\pmb{n}$ 为外法向量.\par
